\begin{summarybox}
\textbf{C011 一页复盘: 记住“一条切线、两个斜率、一个定值”}

\textbf{1. 适用条件}
\[
\frac{x^2}{a^2}+\frac{y^2}{b^2}=1,\quad a>b>0,\quad
P(x_0,y_0)\text{ 在椭圆上},\quad x_0y_0\ne0.
\]

\textbf{2. 三个必须记住的公式}
\[
\frac{x_0x}{a^2}+\frac{y_0y}{b^2}=1,
\qquad k_l=-\frac{b^2x_0}{a^2y_0},
\qquad k_{OP}=\frac{y_0}{x_0}.
\]
\[
\boxed{k_lk_{OP}=-\frac{b^2}{a^2}=e^2-1.}
\]

\textbf{3. 最快证明 (三步)}
\begin{enumerate}
  \item 切线方程读出 $k_l=-b^2x_0/(a^2y_0)$.
  \item 过中心和切点的直线给出 $k_{OP}=y_0/x_0$.
  \item 相乘约去 $x_0,y_0$, 得到 $-b^2/a^2$.
\end{enumerate}

\textbf{4. 最直观记忆}

\textbf{圆: 半径 $\perp$ 切线, 斜率积 $=-1$.}
把圆按 $(x,y)=(au,bv)$ 拉成椭圆,
两条斜率各乘 $b/a$, 乘积自然是 $-(b/a)^2$.
非圆椭圆的一般切点处, $OP$ 并非切线的法线.

\textbf{5. 六种证明路线}

直接用切线方程 (最快); 重根判别式 (无导数);
隐函数求导 (通法); 参数法 (三角消元);
圆到椭圆的仿射拉伸 (几何直觉);
中点弦平方差与割线极限 (联系知识地图).

\textbf{6. 常考联系与变式}
\begin{itemize}
  \item 已知 $k_{OP}$ 求切线斜率: $k_l=-b^2/(a^2k_{OP})$.
  \item 已知 $e$ 求斜率积: $k_lk_{OP}=e^2-1$.
  \item 弦 $AB$ 中点 $M$: 两个斜率存在时,
  $k_{AB}k_{OM}=-b^2/a^2$.
  \item 若中心平移为 $S(h,k)$, 应使用 $SP$ 而非 $OP$.
\end{itemize}

\textbf{7. 最容易丢分的特殊点}

四个轴顶点处: $OP\perp l$, 但至少一条直线没有斜率,
\textbf{斜率积公式不适用}.
非轴点处由于 $-b^2/a^2\ne-1$, 切线与 $OP$ 不垂直.

\textbf{一句口诀}: \textbf{“切线 $x/y$, 径向 $y/x$; 一乘全约掉, 留下负轴比平方.”}
\end{summarybox}
